\section{Classification SVM}

Le classifieur SVM de Matlab est \texttt{fitcsvm}. La commande
\begin{lstlisting}
svm = fitcsvm(X,y);
\end{lstlisting}
renvoie un classifieur SVM entraîné sur les données contenues dans $\X$, avec labels $\yy$. Cette fonction prend par ailleurs plusieurs options (consulter son aide), notamment:
\begin{itemize}[label=-]
	\item \texttt{'BoxConstraint'}: si cette option est fixée à \texttt{inf}, le classifieur utilise une marge dure. Pour toute autre valeur strictement positive, la marge est souple (la valeur par défaut est 1)
	\item \texttt{'KernelFunction'}: le noyau utilisé le cas échéant
	\item \texttt{'KernelScale'}: dans le cas d'un noyau gaussien, ce paramètre correspond à l'écart-type
	\item \texttt{'Verbose'}: affiche ou non les informations liées au classifieur pendant l'optimisation
\end{itemize}

\subsection{Classification}

\begin{enumerate}
	\setcounter{enumi}{7}
	%\item Retrouver ces différentes options dans l'aide de \texttt{fitcsvm}.
	\item Entraîner un classifieur SVM à marge souple et à noyau gaussien sur les données projetées. Attention: on utilisera l'option \texttt{..., 'Standardize', true, ...} pour normaliser les données.
\end{enumerate}

Une fois le classifieur entraîné, les valeurs de labels prédites à partir des données d'entraînement (après ACP) se calculent avec la commande \texttt{predict}, qui peut également renvoyer un score pour chaque label, correspondant à la probabilité que ce label soit le bon.
\begin{lstlisting}
[labels,scores] = predict(svm,X_ACP);
\end{lstlisting}

%Dans le cas de deux classes, disons par exemple 1 et 0, la commande Matlab \texttt{confusionmat} retourne la matrice de confusion associée aux prédictions, c'est-à-dire la matrice \eq{\begin{bmatrix} \mathrm{TP} & \mathrm{FN} \\ \mathrm{FP} & \mathrm{TN} \end{bmatrix}} avec
%\begin{itemize}[label=-]
%	\item TP: nombre de bonnes prédictions de la classe 1 ("true positive")
%	\item FP: nombre de mauvaises prédictions de la classe 1 ("false positive")
%	\item FN: nombre de mauvaises prédictions de la classe 0  ("false negative")
%	\item TN: nombre de bonnes prédictions de la classe 0 ("true negative")
%\end{itemize}
%
%Une mesure possible du taux d'erreur est donné par $\frac{FP+FN}{n}$ où $n$ est le nombre d'échantillons. 
\begin{enumerate}
\setcounter{enumi}{8}
	\item Calculer le taux d'erreur de classification sur les données d'entraînement, c'est-à-dire le ratio entre le nombre de mauvaises classifications et le nombre total d'échantillons.
	\item Réduire la dimension des données de validation avec la même matrice \texttt{U} qui a été utilisée pour les données d'entraînement. Calculer les prédictions et le taux d'erreurs de classification sur le jeu de validation.
\end{enumerate}

\subsection{Réglages des hyperparamètres}

\begin{enumerate}
\setcounter{enumi}{10}
	\item Répéter les codes précédents pour différentes valeurs du paramètre \texttt{'KernelScale'} dans l'intervalle [0.5;20]. Tracer les courbes des erreurs d'apprentissage et de validation correspondant. Quelle est une bonne valeur de ce paramètre?
\begin{lstlisting}
sigma_list = 0.5:0.5:20;
% Insérer votre code ici
for i=1:length(sigma_list)
    % Insérer votre code ici
end
% Insérer votre code ici (plot)
\end{lstlisting}

	\item Répéter la procédure pour deux autres classes (chiffres 6 et 8). Est-il plus difficile de séparer ces classes? La courbe d'erreur de validation est-elle différente de la précédente?

	\item  Une façon courante et pratique de régler les hyperparamètres (comme \texttt{'KernelScale'}) est d'utiliser la méthode de validation croisée. Cette méthode n'utilise que les données d'apprentissage, en les sous-divisant d'une manière astucieuse en sous-ensembles d'entraînement et de validation. On peut optimiser le paramètre \texttt{'KernelScale'} de la façon suivante:
\begin{lstlisting}
svm_opt = fitcsvm(X,y,'Verbose',1, ...
	'Standardize',true, ...
	'KernelFunction','gaussian', ...
	'KernelScale',1, ...
	'OptimizeHyperparameters','KernelScale');
sigma_opt = s_opt.KernelParameters.Scale;
\end{lstlisting}

	Comparer les résultats (paramètre optimal, taux d'erreur) avec l'erreur de validation obtenue à l'étape précédente.
	
	\end{enumerate}
	
	\subsection{Pré-traitements additionnels}
	
	\begin{enumerate}
	\setcounter{enumi}{13}
		\item Prenez deux photos (noir et blanc, carré) avec votre smartphone, représentant les deux classes (6 et 8), ou utilisez les photos fournies dans les fichiers du tp. Affichez-les dans Matlab (NB: une image peut-être lue depuis un fichier avec la commande \texttt{imread}). Classifiez-les après les étapes de pré-traitement suivantes:
	\begin{itemize}[label=-]
		\item Si l'image est en couleurs, conversion en niveaux de gris (en prenant par exemple la moyenne des canaux RGB)
		\item Recadrage (si nécessaire)
		\item Négatif (blanc sur fond noir)
		\item Correction du contraste -- il faut normaliser les données de telle façon que les valeurs minimales et maximales deviennent 0 et 255: $f_{i,j} = \frac{f_{i,j}-f_{\min}}{f_{\max}-f_{\min}}\cdot 255$
		\item Redimensionnement à l'aide de la fonction \texttt{imresize}.
	\end{itemize}
\end{enumerate}